% Truncation boundary of the complete dodecaphase route and exact
% delimitation of the finite analytic resolvent.
\paragraph{Exact dependence on the truncation boundary.}
\label{ampl-alfa-frontera-exacta}
The regional coordinates and the dodecaphase register arise from the
common APP--TRIT--TPK generation. Once produced, their positional
composition has an explicit dependence on readout depth. Let $B=1000$,
\[
 Z_j=P_j+E_j-\Phi_j-K_j^{\mathrm{per}},\qquad
 y_N=\sum_{j=1}^{N}Z_jB^{-j},\qquad
 R_{N+1}(Z)=\sum_{r\geq0}Z_{N+1+r}B^{-r-1}.
\]
The complete series in \eqref{eq:alpha-serie} satisfies exactly
\[
 \alpha_A=y_N+B^{-N}R_{N+1}(Z),\qquad
 -2<R_{N+1}(Z)<2.
\]
This identity combines the contribution of the window and that of its
prolongation, both expressed in the same coordinates.

\begin{proposition}[Boundary transport and canonical cylinder]
\label{ampl-alfa-prop-frontera}
Let $A_j^{(N,t)}$ be the output of finite positional division with
$c_{N+1}=t\in\mathcal C=\{-2,-1,0,1,2\}$, and let
$a_N^{(t)}=\sum_{j=1}^{N}A_j^{(N,t)}B^{-j}$. For the records in the
domain under consideration, one has
\begin{equation}
 \alpha_A-a_N^{(t)}
 =B^{-N}\bigl(R_{N+1}(Z)-t\bigr),
 \qquad
 |\alpha_A-a_N^{(t)}|<4B^{-N}.
 \label{ampl-alfa-eq-frontera}
\end{equation}
The quotient of the complete prolongation is
$c_{N+1}^{\infty}=\lfloor R_{N+1}(Z)\rfloor$. With this boundary,
\[
 0\leq\alpha_A-a_N^{(c_{N+1}^{\infty})}<B^{-N},
\]
so that the resulting window is the canonical prefix of order $N$.
\end{proposition}
\begin{proof}
Telescoping \eqref{eq:alpha-identidad-local} gives
$a_N^{(t)}=y_N+tB^{-N}-c_1$. Since $Z_1=7$, every incoming state
belongs to $\mathcal C$, and the first dividend lies between $5$ and $9$,
one has $c_1=0$. Substituting this equality into the decomposition of the
series proves \eqref{ampl-alfa-eq-frontera}. The bound follows from
$R_{N+1}(Z)\in(-2,2)$ and $t\in\mathcal C$.

For the complete canonical normalization,
$c_{N+1}^{\infty}=R_{N+1}(Z)-R_{N+1}(A)$ and
$R_{N+1}(A)\in[0,1)$. Integrality of the quotient implies
$c_{N+1}^{\infty}=\lfloor R_{N+1}(Z)\rfloor$. The error formula
then becomes
$\alpha_A-a_N^{(c_{N+1}^{\infty})}=B^{-N}R_{N+1}(A)$, which proves
membership in the canonical cylinder.
\end{proof}

The proposition identifies the meaning of the right boundary: it is the
integer part of the signed tail expressed on the scale of the next
position. Computation over the five boundaries preserves this information
as a finite collection of possibilities until the blocks to the left
stabilize. The quantitative contraction arises from depth and tail
transport; its object is the normalization of the coordinates already generated.

\paragraph{Graded domain and status of the analytic representation.}
\label{ampl-alfa-dominio-resolvente}
The analytic composition displayed in
\S\ref{subsec:alpha-resolvente} acts on
$\operatorname{span}\{e_7,e_8,e_9\}$ and satisfies $Ne_9=0$. Its
evaluation produces exactly $T_{7:9}$; degree $9$ is the last
component of that domain. The identity $N^3=0$ explains why the
resolvent is expressed by three summands. The local resolvent does not by
itself determine coefficients of higher degree.
Section~\ref{subsec:alpha-limite-dos-vias} proves this nonuniqueness of the
bare jet and separates two operations: the dodecaphase representation produces the
coordinate \(\alpha_{\mathrm{HMT}}\) from the state records; subsequently,
the analytic representation consumes the same precoordinate and generates, through an
explicit recurrence, all coefficients of its analytic image. No
conventional root is consulted, nor is the analytic representation promoted to an
independent generator. The compatibility squares and nested cylinders
prove equality of both representations at every depth. Local
nilpotence, dodecaphase representation, and analytic representation
thus perform distinct functions: the first certifies the finite result
of order nine; the second determines \(\alpha_{\mathrm{HMT}}\); the third
represents that same coordinate at arbitrary depth and verifies its compatibility.
